\section{Untersuchung des Weiterentwicklungsbedarfs}
etwas empirisch etwas nicht

in 4.3 was gemacht werden muss.
hier was gemacht werden kann/soll.

\subsection{Ablauf der Untersuchung und Methodik}

\subsection{Risikoanalyse bei Nicht-Realisierung des Projekts \enquote{Mailbroadcast Stufe 3}}
\subsubsection{Anforderungen an eine Broadcast-Lösung zur Einhaltung einer rechtssicheren Unternehmens-, Krisen- und Notfallkommunikation}
\subsubsection{Anforderungserfüllung der Anwendung Mailbroadcast im Ist- und Soll-Zustand}
darauf eingehen, inwiefern der Ist-Zustand von MBC die Anforderungen aus dem vorherigen Kapitel erfüllt

\subsection{Befragung der Endanwender}


\subsection{Qualität bereitgestellter Personaldaten}
\subsubsection{Methoden zur Messung der Datenqualität}
\subsubsection{Datenanalyse und Auswertung}
\subsubsection{Beurteilung}

\subsection{Erforderliche Maßnahmen zur Nutzung von Mailbroadcast in der Notfallkommunikation}
\label{subsec:mbc_notfall}
%\subsection{Priorisierung von Notfallmails}
%\subsubsection{Notwendig?}
%\subsubsection{Befragung}
%\subsubsection{Analyse verschiedener Szenarien (vll. vergangener Krisenfälle)}
%\subsubsection{Auswirkung auf die gesamte Kommunikation}

\subsection{Synthese}